\scanpage{52}
\uline{Lemme} Soient $\deleted{\unsure{k}}K$ un corps, $G$ un groupe algébrique sur $\deleted{k}K$,
$M=(M(n))$ un groupe p-divisible ($p$ un \unsure{nb|nt} premier),
$\varphi:M\unsure{\to|-}G$ un hom de groupes. \inserted{On suppose que $\mathrm{car}\,K=p$, ou \deleted{que} $K$ parfait, et que $G$ ne contient pas de ss-groupe isomorphe à $G_a$.} Alors \inserted{\deleted{a) Si $\mathrm{car}\,K\ne p$,}} il existe
un sous-schéma en groupes $H$ de $G$, extension
d'une VA par un tore, tel que $\varphi$ se factorise
par $H$. \deleted{b) Si $\mathrm{car}\, K=$}

\uline{Dém} On a $\varphi=(\varphi(n))$, avec
\[
\varphi(n):M(n)\longrightarrow G
\]
Soit $Z_n=\mathrm{Centr}\,\varphi(n)$. Les $Z_n$ sont des sous-schémas en groupes formant une suite décroissante,
donc stationnaire, soit $Z$ la valeur commune.
Évidemment $\varphi$ se factorise par \deleted{$\unsure{G}$} $Z$, et
même par $\mathrm{Centr}(Z)$. Cela nous ramène au
cas où $G$ est commutatif. Evidemment
\deleted{\unsure{on a}} l'hom composé $M\unsure{-|\to}G-G/G^0$ est nul, donc
on peut supposer $G$ connexe.

\uline{1\textsuperscript{er} Cas} $p=\mathrm{car}. K$. On a évidemment, pour tt
entier $\unsure{\nu|v}\geq 0$, que $\varphi$ se factorise par $p^\nu G$.
Or pour $\nu$ grand, on sait \deleted{$p^\nu G$} que $p^\nu G$ est
une extension d\unsure{\textsuperscript{n}|'} une VA par un tore (à dégager
en lemme, pour tt groupe algébrique commutatif
sur corps de car. $p>0$).

\uline{2\textsuperscript{ème} Cas} \deleted{$\unsure{p\ne\mathrm{car}\,K}$. Les $M(n)$ sont alors étales.} \inserted{\deleted{\unsure{car|ou}} $K$ parfait, $G$ n'a pas de ss-grpe \unsure{isom :|isom. à} $G_a$.}
Alors on sait que $G$ est ext. d'une VA. par un tore !
% archive sheet number: 51
